%% =====================================================================
%%  Annexe — Mémento de grammaire (toutes les règles du tome en tableaux)
%% =====================================================================
\annexe[Memento]{Mémento de grammaire}{Toutes les règles du tome, en tableaux}

\noindent{\fontsize{10.4}{14}\selectfont Ce mémento rassemble toute la grammaire du tome 1. Pour chaque
point, l'unité indiquée entre parenthèses vous renvoie aux explications complètes et aux exercices.\par}

%% ---------------------------------------------------------------------
\otIntertitre[otNavy]{1 · Les pronoms et les possessifs}
\begin{center}\otTableau
\begin{tabular}{@{}l@{\hspace{10mm}}>{\enOn}c<{\enOff}@{\hspace{10mm}}>{\enOn}c<{\enOff}@{\hspace{10mm}}>{\enOn}c<{\enOff}@{\hspace{10mm}}l@{}}
 & \otTH{Sujet (u. 1)} & \otTH{Complément (u. 15)} & \otTH{Possessif (u. 4)} & \otTH{Français}\\\midrule
1\textsuperscript{re} pers. sing. & I    & me   & my    & je, me, mon\\
2\textsuperscript{e} pers.        & you  & you  & your  & tu / vous, te, ton\\
3\textsuperscript{e} masc.        & he   & him  & his   & il, le, son\\
3\textsuperscript{e} fém.         & she  & her  & her   & elle, la, son\\
3\textsuperscript{e} neutre       & it   & it   & its   & il / elle (chose), le, son\\
1\textsuperscript{re} pers. plur. & we   & us   & our   & nous, notre\\
3\textsuperscript{e} pers. plur.  & they & them & their & ils / elles, les, leur\\\bottomrule
\end{tabular}\end{center}
{\fontsize{9.8}{12.6}\selectfont Le possessif s'accorde avec le possesseur : \en{his mother} (la mère de
Jack), \en{her father} (le père d'Emma). Le génitif : \en{Léa's brother}, \en{my parents' house} (u.~4).\par}

%% ---------------------------------------------------------------------
\otIntertitre[otNavy]{2 · \en{be} au présent (unités 1 et 2)}
\begin{center}\otTableau
\begin{tabular}{@{}>{\enOn}l<{\enOff}@{\hspace{8mm}}>{\enOn}l<{\enOff}@{\hspace{8mm}}>{\enOn}l<{\enOff}@{\hspace{8mm}}>{\enOn}l<{\enOff}@{}}
\otTH{Affirmation} & \otTH{Négation} & \otTH{Question} & \otTH{Réponses courtes}\\\midrule
I'm (I am)          & I'm not            & Am I…?     & Yes, I am. / No, I'm not.\\
you're (you are)    & you aren't         & Are you…?  & Yes, you are. / No, you aren't.\\
he's, she's, it's   & he isn't           & Is he…?    & Yes, he is. / No, he isn't.\\
we're, they're      & we aren't          & Are they…? & Yes, they are. / No, they aren't.\\\bottomrule
\end{tabular}\end{center}
{\fontsize{9.8}{12.6}\selectfont L'âge se dit avec \en{be} : \en{I'm twenty-eight} (jamais \en{I have
twenty-eight}). Pas de contraction dans une réponse courte affirmative : \en{Yes, I am}.\par}

%% ---------------------------------------------------------------------
\begin{paire}
\begin{regle}[Articles et pluriel · unité 5]
\en{a} + son consonne : \en{a pen}, \en{a university}\par
\en{an} + son voyelle : \en{an apple}, \en{an hour}\par
\en{the} : chose précise, connue\par\vspace{1mm}
Pluriel : \en{-s} · \en{-es} (\en{boxes}, \en{watches}) · \en{-ies} (\en{cities})\par
Irréguliers : \en{men}, \en{women}, \en{children}, \en{people}\par\vspace{1mm}
\en{this} / \en{these} (ici) · \en{that} / \en{those} (là-bas)
\end{regle}
\begin{regle}[Les mots interrogatifs · unités 3 et 20]
\en{What?} quoi · \en{Where?} où · \en{Who?} qui\par
\en{When?} quand · \en{Why?} pourquoi · \en{How?} comment\par
\en{How old?} quel âge · \en{What time?} à quelle heure\par
\en{How much?} combien (prix, indénombrable)\par
\en{How many?} combien (pluriel)\par
\en{How long?} combien de temps
\end{regle}
\end{paire}

%% ---------------------------------------------------------------------
\otIntertitre[otNavy]{3 · Le présent simple : les habitudes (unités 6 et 7)}
\begin{center}\otTableau
\begin{tabular}{@{}>{\enOn}l<{\enOff}@{\hspace{8mm}}>{\enOn}l<{\enOff}@{\hspace{8mm}}>{\enOn}l<{\enOff}@{\hspace{8mm}}>{\enOn}l<{\enOff}@{}}
\otTH{Affirmation} & \otTH{Négation} & \otTH{Question} & \otTH{Réponses courtes}\\\midrule
I / you / we / they work & I don't work     & Do you work?    & Yes, I do. / No, I don't.\\
he / she / it work\textbf{s} & she doesn't work & Does she work? & Yes, she does. / No, she doesn't.\\\bottomrule
\end{tabular}\end{center}
{\fontsize{9.8}{12.6}\selectfont Un seul \en{-s} par phrase : \en{she works}, mais \en{she doesn't work},
\en{Does she work?} Orthographe : \en{watches}, \en{goes}, \en{studies}, \en{has}. Mots repères :
\en{every day}, \en{usually}, \en{always}, \en{on Mondays}. Les goûts : \en{like}, \en{love}, \en{hate} +
nom ou verbe en \en{-ing} (\en{I love cooking}).\par}

%% ---------------------------------------------------------------------
\begin{paire}
\begin{regle}[\en{there is} et le lieu · unités 8 et 9]
\en{There's a kitchen.} · \en{There are two rooms.}\par
\en{Is there…? — Yes, there is.} · \en{There aren't any…}\par\vspace{1mm}
\en{in} · \en{on} · \en{under} · \en{behind} · \en{in front of} · \en{next to} · \en{between} ·
\en{opposite} · \en{on the corner of}
\end{regle}
\begin{regle}[L'impératif · unité 9]
\en{Turn left.} · \en{Go straight ahead.} · \en{Take the subway.}\par
Négation : \en{Don't take the bus.} · \en{Don't be late!}\par\vspace{1mm}
Proposer : \en{Let's go!} (u.~15)
\end{regle}
\end{paire}

\begin{paire}
\begin{regle}[Quantités · unité 10]
\en{How much is it? — It's twenty dollars.}\par
\en{How much water?} (indénombrable)\par
\en{How many apples?} (dénombrable)\par\vspace{1mm}
\en{some} : affirmation (\en{I need some milk})\par
\en{any} : négation et question (\en{Are there any eggs?})\par
\en{a lot of} : beaucoup de
\end{regle}
\begin{regle}[Demander poliment · unité 11]
\en{Can I have a coffee, please?}\par
\en{Could I have the fish?} (plus poli)\par
\en{I'd like a salad.} (= \en{I would like})\par
\en{I'll have the burger.}\par\vspace{1mm}
\en{Would you like something to drink?}\par
\en{Would you like to come?} : invitation (u.~15)
\end{regle}
\end{paire}

%% ---------------------------------------------------------------------
\otIntertitre[otNavy]{4 · \en{can} : capacité, permission (unités 11 et 12)}
\begin{center}\otTableau
\begin{tabular}{@{}>{\enOn}l<{\enOff}@{\hspace{8mm}}>{\enOn}l<{\enOff}@{\hspace{8mm}}>{\enOn}l<{\enOff}@{\hspace{8mm}}>{\enOn}l<{\enOff}@{}}
\otTH{Affirmation} & \otTH{Négation} & \otTH{Question} & \otTH{Réponses courtes}\\\midrule
I / she / they can swim & He can't (cannot) ski & Can you drive? & Yes, I can. / No, I can't.\\\bottomrule
\end{tabular}\end{center}
{\fontsize{9.8}{12.6}\selectfont Ni \en{-s}, ni \en{to}, ni \en{do} : \en{she can drive}. Le niveau :
\en{very well}, \en{quite well}, \en{not very well}, \en{not at all}. \en{good} est un adjectif (\en{a good
swimmer}), \en{well} un adverbe (\en{she swims well}).\par}

%% ---------------------------------------------------------------------
\otIntertitre[otNavy]{5 · Le présent en \en{-ing} : maintenant (unité 13)}
\begin{center}\otTableau
\begin{tabular}{@{}>{\enOn}l<{\enOff}@{\hspace{8mm}}>{\enOn}l<{\enOff}@{\hspace{8mm}}>{\enOn}l<{\enOff}@{\hspace{8mm}}>{\enOn}l<{\enOff}@{}}
\otTH{Affirmation} & \otTH{Négation} & \otTH{Question} & \otTH{Réponses courtes}\\\midrule
I'm working.     & I'm not working.   & Are you working? & Yes, I am. / No, I'm not.\\
She's working.   & She isn't working. & Is she working?  & Yes, she is. / No, she isn't.\\\bottomrule
\end{tabular}\end{center}
{\fontsize{9.8}{12.6}\selectfont Orthographe : \en{make} → \en{making} ; \en{sit} → \en{sitting}.
Mots repères : \en{now}, \en{right now}, \en{at the moment}, \en{Look!}, \en{Listen!}. Les verbes d'état
restent au présent simple : \en{I want}, \en{I like}, \en{I know}, \en{I understand}.\par}

%% ---------------------------------------------------------------------
\begin{paire}
\begin{regle}[Dates et prépositions de temps · unité 14]
Ordinaux : \en{first}, \en{second}, \en{third}, \en{fourth}… \en{fifth}, \en{twelfth},
\en{twentieth}, \en{twenty-first}\par\vspace{1mm}
Date américaine : \en{March 15th} (dit \en{March fifteenth}) ; écrite \en{3/15}\par\vspace{1mm}
\en{at} + heure · \en{on} + jour, date · \en{in} + mois, saison, année\par
\en{at night} · \en{in the morning}
\end{regle}
\begin{regle}[\en{be going to} · unité 16]
\en{I'm going to travel.}\par
\en{She isn't going to fly.}\par
\en{Are you going to stay? — Yes, I am.}\par\vspace{1mm}
Projet décidé ou prévision évidente (\en{It's going to rain}).\par
Repères : \en{tomorrow}, \en{tonight}, \en{next week} (sans \en{the}), \en{in two weeks}, \en{soon}
\end{regle}
\end{paire}

%% ---------------------------------------------------------------------
\otIntertitre[otNavy]{6 · Le passé : \en{was}, \en{were} et le prétérit (unités 17 à 19)}
\begin{center}\otTableau
\begin{tabular}{@{}l@{\hspace{6mm}}>{\enOn}l<{\enOff}@{\hspace{6mm}}>{\enOn}l<{\enOff}@{\hspace{6mm}}>{\enOn}l<{\enOff}@{}}
 & \otTH{\en{be}} & \otTH{Verbe régulier} & \otTH{Verbe irrégulier}\\\midrule
affirmation      & I was · they were           & I watched         & I went\\
négation         & I wasn't · they weren't     & I didn't watch    & I didn't go\\
question         & Was she…? · Were you…?      & Did you watch…?   & Did you go…?\\
réponses courtes & Yes, I was. · No, I wasn't. & Yes, I did.       & No, I didn't.\\\bottomrule
\end{tabular}\end{center}
{\fontsize{9.8}{12.6}\selectfont Orthographe du \en{-ed} : \en{loved}, \en{tried}, \en{stopped},
\en{played}. Prononciation : \phon{t} (\en{walked}), \phon{d} (\en{played}), \phon{ɪd} (\en{visited}).
\en{there was} / \en{there were} ; \en{two days ago} ; \en{I was born in 1998}. Après \en{did} et
\en{didn't}, toujours la base verbale. Liste des irréguliers page~\pageref{annexe-Verbes}.\par}

%% ---------------------------------------------------------------------
\otIntertitre[otNavy]{7 · Récapitulatif : choisir le temps et l'auxiliaire (unité 20)}
\begin{center}\otTableau
\begin{tabular}{@{}l@{\hspace{6mm}}>{\enOn}l<{\enOff}@{\hspace{6mm}}>{\enOn}l<{\enOff}@{\hspace{6mm}}>{\enOn}l<{\enOff}@{}}
\otTH{Pour parler\ldots} & \otTH{On dit\ldots} & \otTH{Question} & \otTH{Mots repères}\\\midrule
d'une habitude        & I work.                 & Do you work?               & every day, usually\\
d'une action en cours & I'm working.            & Are you working?           & now, Look!\\
d'un projet           & I'm going to work.      & Are you going to work?     & tomorrow, next week\\
d'une capacité        & I can work.             & Can you work?              & —\\
du passé (\en{be})    & I was at work.          & Were you at work?          & yesterday, ago\\
du passé (autres verbes) & I worked. · I went.  & Did you work? · Did you go? & yesterday, last week\\\bottomrule
\end{tabular}\end{center}
